\paragraph{The integral determinant at ramified primes.}
The image of an eigenprojection and its integral eigenkernel need not
coincide. Their exact discrepancy follows from one cokernel calculation.

\begin{lemma}\label{geo-ramified-resultant-index}
Let $R$ be a discrete valuation ring with fraction field $F$, and let
$a_i\in R$, $i\in I$, be distinct, with $I$ finite and nonempty.
Put $B=R[t]/\prod_i(t-a_i)$ and $H=B^n$.
Suppose $L$ is a $B$-linear quotient of $H$, finite free over $R$,
and write $K$ for its kernel. Let $e_i$ be the rational coefficient
idempotents, $L_i=e_iL_F$, $d_i=\dim_F L_i$,
$L^i=L\cap L_i$ and $Q_i=e_iH$.
Then the rational determinant isomorphism gives
\begin{equation}\label{geo-ramified-resultant-index-eq}
 \det(e_iK)\otimes\det L^i
   =\left(\prod_{j\ne i}(a_i-a_j)^{d_j}\right)\det Q_i
\end{equation}
as rank-one $R$-lattices. Empty determinants and products equal one.
\end{lemma}
\begin{proof}
Write $T$ for multiplication by $t$. Since
$H/(T-a_i)H=Q_i$, right exactness gives
\[
 Q_i/e_iK\simeq L/(T-a_i)L.
\]
Set $L_0=L\cap\bigoplus_{j\ne i}L_j$ and $P_0=(1-e_i)L$.
Projection identifies
\[
 P_0/L_0\simeq e_iL/L^i,
 \qquad
 \operatorname{tors}(L/(T-a_i)L)=L_0/(T-a_i)P_0.
\]
The exact sequence
\[
 0\longrightarrow L_0/(T-a_i)P_0
 \longrightarrow P_0/(T-a_i)P_0
 \longrightarrow P_0/L_0\longrightarrow0
\]
therefore shows that the sum of these two lengths is
$\sum_{j\ne i}d_jv(a_i-a_j)$: it is the valuation of the
determinant of $T-a_i$ on $\bigoplus_{j\ne i}L_j$.
Here $v$ takes value one on a uniformizer.

The kernel $K_0$ of $Q_i\to e_iL$ is saturated, and
$K_0/e_iK$ is the torsion module just displayed.
A basis of $e_iK$, together with lifts of a basis of $L^i$,
has index
$\operatorname{length}(K_0/e_iK)+
 \operatorname{length}(e_iL/L^i)$ in $Q_i$.
This proves \eqref{geo-ramified-resultant-index-eq}.
The order of the determinant factors changes only a unit sign.
\end{proof}

Let $E$ be a number field with involution $c$, and let $K$ contain
its normal closure. Fix a signature $d_\tau$, indexed by the
embeddings $\tau:E\hookrightarrow K$. At a finite place $v$ of $K$,
write $R=\mathcal O_{K_v}$ and let $S_{\sigma,v}$ consist of the
embeddings with the same residue character as $\sigma$.
The corresponding factor of $\mathcal O_E\otimes R$ is
\[
 B_{\sigma,v}=R[t]/\prod_{\tau\in S_{\sigma,v}}(t-a_\tau).
\]
Indeed, the largest unramified subfield of the relevant completion of
$E$ splits over $R$, and a uniformizer generates its remaining totally
ramified integer ring. This holds also at $2$ and with wild ramification.
\cite[Proposition 7.50, Corollary 7.51, Proposition 7.55 and
Remark 7.56]{milne-ant}
Define an ideal of $\mathcal O_K$ by
\begin{equation}\label{geo-signature-correction-ideal}
 \mathfrak c_\sigma R
   =\left(\prod_{\substack{\tau\in S_{\sigma,v}\\\tau\ne\sigma}}
                   (a_\sigma-a_\tau)^{d_\tau}\right)R.
\end{equation}
Each local ideal is independent of the generator: two generators
are polynomials in one another, so their character differences divide
one another. The ideal is the unit ideal away from the discriminant
of $E$, and its formation commutes with finite extensions of splitting
valuation rings.

\begin{proposition}\label{geo-ramified-hodge-determinant}
Let $\mathcal Y/\mathcal O_K$ be a flat arithmetic model carrying
its prescribed tangent bundle $\mathcal L$ and an exact Hodge sequence
\[
 0\longrightarrow\mathcal L^\vee\longrightarrow\mathcal H
   \longrightarrow\mathcal L\longrightarrow0.
\]
Assume all terms are locally free over $\mathcal O_{\mathcal Y}$,
the sequence is $\mathcal O_E$-equivariant, and the kernel action is
$a\lambda=\lambda\circ c(a)$.
Assume $\mathcal H$ is locally free of rank $n$ over
$\mathcal O_E\otimes\mathcal O_{\mathcal Y}$ and the generic
signature of $\mathcal L$ is $d_\tau$, constant on the selected
component. Then $d_\tau+d_{\tau c}=n$, and
\begin{equation}\label{geo-actual-integral-determinant}
 \mathcal Q_\sigma
   =\mathcal H\otimes_{\mathcal O_E\otimes\mathcal O_{\mathcal Y},\sigma}
                      \mathcal O_{\mathcal Y},
 \qquad
 \mathcal M_\sigma
   =\det\mathcal Q_\sigma\otimes
                         \operatorname{pr}^*\mathfrak c_\sigma
\end{equation}
are a vector bundle and an invertible sheaf, respectively.
On the generic fibre $\mathcal M_\sigma$ is the conjugate tangent
determinant. At every integral point, also after finite ramified
extension, its pullback is exactly
\[
 \det L^\sigma\otimes(\det L^{\sigma c})^{-1},
\]
where these are the actual saturated tangent eigenlattices.
On every special-fibre curve its degree equals that of the character
determinant $\det\mathcal Q_\sigma$.
\end{proposition}
\begin{proof}
Work first over a splitting complete valuation ring $R$.
The rational Hodge sequence identifies
$\det(H_F)_\sigma$ with
$\det(L_F)_\sigma\otimes\det(L_F)_{\sigma c}^{-1}$.
Integrally, restriction gives
\[
 e_\sigma(L^\vee)=(L^{\sigma c})^\vee:
\]
the lattice $L^{\sigma c}$ is saturated in $L$, so every integral
functional on it extends to $L$. The appropriate factor of $H$ is
free over the local ring $B_{\sigma,v}$.
Lemma~\ref{geo-ramified-resultant-index} gives
\[
 \det L^\sigma\otimes(\det L^{\sigma c})^{-1}
      =(\mathfrak c_\sigma R)\det Q_\sigma.
\]
This proves the claimed integral identification. Completion is
faithfully flat; coefficient eigenkernels commute with flat extension,
as kernels of finitely many action equations. Formula
\eqref{geo-signature-correction-ideal} commutes with finite extension,
so the same conclusion holds for every stated integral point.
The base ideal is an invertible sheaf and restricts to a constant
one-dimensional line on each special-fibre curve. Its degree there
is zero.
\end{proof}

For the prescribed global family, the Hodge extension in the proposition
is constructed directly from the original rigidified biextension
$\mathcal P$. This first-order part of the construction of
Theorem~\ref{geo-biextension-conjugate-filtration} works in mixed
characteristic. If $G_1$ is the first neighbourhood of zero and
$q:G\times G_1\to G$, then
\[
 0\longrightarrow\operatorname{pr}_G^*\omega_G
 \longrightarrow q_*\mathcal P\longrightarrow\mathcal O_G
 \longrightarrow0.
\]
The inverse image of $1$ is an additive $\omega_G$-torsor. Its group
law comes from the first biextension law, and its Lie algebra is
$\mathcal H$. Global adjunction, compatible with addition,
multiplication, rigidification and symmetry, defines its coefficient
action by pullback along $c(a)$ on $G_1$ followed by inverse adjunction.
All these operations commute with base change and glue on the actual
charts. Thus the kernel is the exact integral conjugate dual.

On the actual integral model, whose group has abelian generic fibre,
the global adjunction isomorphisms admit a unique coherent normalization.
Every regular function on $G^m$ equals its zero-section value: this holds
on the proper geometrically connected generic fibre, and flatness makes
that fibre schematically dense. Divide each adjunction by its origin
unit. Any two isomorphisms of the same line bundles with that normalized
origin value are consequently equal. For addition the common source
and target are $\mathcal P((a+b)x,y)$ and $\mathcal P(x,c(a+b)y)$,
using the two biextension laws. For multiplication the intermediate
line is $\mathcal P(bx,c(a)y)$ and the final line is
$\mathcal P(x,c(b)c(a)y)$. The biextension laws similarly compare
splitting $x+x'$ or $y+y'$. Symmetry compares adjunction followed by
exchange with exchange followed by inverse adjunction for $c(a)$.
All canonical identifications preserve the rigidification, so these
diagrams commute, as do the unit and zero diagrams. This supplies the
required coherence from the original global cubical data.

Under the stated self-dual quotient hypotheses, the proper and toric
comparisons make every geometric fibre projective over its coefficient
algebra. To recover the common rank $n$, use an integral basis $b_j$
of $\mathcal O_E$: the polynomial
$\det(T-\sum_jX_jb_j\mid\mathcal H)$ is the $n$-th power of the
regular coefficient polynomial generically, hence everywhere by
flatness. At a geometric special point its distinct character factors
have multiplicities $m_\chi r_\chi$, including the full local Artin
length $m_\chi$. Comparison gives $r_\chi=n$.
This freeness descends to the closed residue field: on each Artin local
factor, lift a residue-field basis and compare dimensions after extending
the field. At a stalk with local ring $A$, lift a coefficient basis of
the closed fibre. The resulting
map from $(\mathcal O_E\otimes A)^n$ to $\mathcal H_A$ is a map of
free $A$-modules of the same rank, invertible modulo the maximal
ideal, and hence invertible. This proves coefficient local freeness
on the same model. It requires the stated global biextension and
adjunction isomorphisms; compatible charts only on test curves do not assert a
global family.

The proposition constructs the ramified integral line without assuming
that the individual tangent eigensheaves are locally free over the
coefficient algebra. If the actual curvewise Hodge and conjugate
filtrations satisfy Proposition~\ref{geo-frobenius-determinant-torsion},
the character determinant has degree zero. If the model is normal and
proper with nonempty connected special fibres,
Theorem~\ref{geo-vertical-line-rigidity} consequently makes a determinant
section generating the whole generic fibre have point-independent finite norms on
these exact tangent eigenlattices.
The arithmetic model and its prescribed pointwise period data remain
required. Propositions~\ref{geo-cm-period-point-descent}
and~\ref{geo-cm-integral-map-extension} below supply the exact CM maps
from those data and moduli point separation.
